\newpage

\section{Arbres Binomiaux et Évaluation d'Options}

\subsection{Le Modèle à une Période (L'Univers à Deux États)}

\begin{intuitionbox}[Le concept de l'arbre]
Pour évaluer une option classique (vanille), une méthode redoutablement efficace consiste à modéliser les mouvements de l'actif sous-jacent à l'aide d'un arbre. Nous commençons par une situation hautement stylisée (un seul pas de temps, deux états futurs possibles) afin de comprendre la mécanique fondamentale de l'absence d'arbitrage, avant de généraliser.
\end{intuitionbox}

\begin{definitionbox}[Le cadre du modèle binomial simple]
Considérons un univers simplifié avec les paramètres suivants :
\begin{itemize}[label=$\rightarrow$]
    \item \textbf{Taux d'intérêt} : Pour simplifier, nous supposons qu'il est nul ($r = 0$).
    \item \textbf{Actif sous-jacent ($S$)} : Son prix actuel est $S_0 = 100$. Demain, il ne peut prendre que deux valeurs :
    \begin{itemize}[label=$\cdot$]
        \item État "Up" (Hausse) : $S_u = 110$
        \item État "Down" (Baisse) : $S_d = 90$
    \end{itemize}
    \item \textbf{Option d'achat (Call)} : Prix d'exercice $K = 100$.
    \begin{itemize}[label=$\cdot$]
        \item Flux en cas de Hausse : $\max(110 - 100, 0) = 10$
        \item Flux en cas de Baisse : $\max(90 - 100, 0) = 0$
    \end{itemize}
\end{itemize}
\end{definitionbox}

\subsection{Tarification par Couverture (Hedging)}

\begin{intuitionbox}[L'élimination de l'incertitude]
En tant que vendeur de cette option, nous sommes exposés à un risque : perdre 10 si l'état "Up" se réalise, et ne rien perdre dans l'état "Down". La stratégie de couverture consiste à construire un portefeuille aujourd'hui (composé de l'actif sous-jacent et de l'option) dont la valeur de demain sera strictement identique, quel que soit l'état du monde qui se réalise.
\end{intuitionbox}

\begin{proofbox}[Détermination du prix sans arbitrage par la couverture]
Sans connaître l'avenir, nous achetons une quantité fixe $\delta$ d'actions aujourd'hui. Soit $Opt_0$ la valeur de l'option aujourd'hui.
\begin{itemize}[label=$\rightarrow$]
    \item \textbf{Valeur du portefeuille aujourd'hui} : $\Pi_0 = 100\delta - Opt_0$
\end{itemize}
Regardons la valeur de ce portefeuille demain dans les deux états :
\begin{itemize}[label=$\rightarrow$]
    \item Dans l'état "Up" : $\Pi_u = 110\delta - 10$
    \item Dans l'état "Down" : $\Pi_d = 90\delta$
\end{itemize}

\textbf{Étape 1 : Couverture du risque (Delta)}\\
Pour éliminer toute incertitude, le portefeuille doit valoir la même chose dans les deux états ($\Pi_u = \Pi_d$) :
\begin{equation*}
    110\delta - 10 = 90\delta \implies 20\delta = 10 \implies \delta = \frac{1}{2}
\end{equation*}
En détenant une demi-action, notre portefeuille vaudra demain exactement $90 \times (0,5) = 45$, quoi qu'il arrive.

\textbf{Étape 2 : Application du principe de non-arbitrage}\\
Puisque le taux d'intérêt est nul ($r=0$) et que notre portefeuille est sans risque, sa valeur aujourd'hui doit être exactement égale à sa valeur de demain :
\begin{equation*}
    \Pi_0 = 45 \implies 100(0,5) - Opt_0 = 45 \implies 50 - Opt_0 = 45 \implies Opt_0 = 5
\end{equation*}
Le prix unique et sans arbitrage de l'option est de $5$.
\end{proofbox}

\begin{remarquebox}[L'absence des probabilités]
Le point le plus remarquable de ce raisonnement est que les probabilités n'interviennent nulle part. Peu importe qu'il y ait 90\% ou 10\% de chances que l'action monte : l'argument de couverture reste parfaitement valide. La seule condition imposée aux probabilités réelles est qu'elles ne soient ni $0$ ni $1$ (car une probabilité de $1$ de voir l'action monter plus vite que le taux sans risque créerait une opportunité d'arbitrage immédiate).
\end{remarquebox}

\newpage

\subsection{L'Évaluation Risque-Neutre}

\begin{intuitionbox}[L'approche probabiliste]
Que se passe-t-il si nous essayons d'évaluer l'option en utilisant une espérance mathématique ? Pour que cette approche donne le même résultat que notre méthode de couverture (qui est infaillible), nous devons inventer un monde théorique où les investisseurs sont neutres au risque. Dans ce monde, le rendement attendu de l'action est égal au taux sans risque.
\end{intuitionbox}

\begin{theorembox}[Probabilité Risque-Neutre]
Pour retrouver le prix d'arbitrage, on définit une probabilité fictive $p$ (pour l'état "Up") et $1-p$ (pour l'état "Down"), telle que l'espérance de la valeur future de l'action soit exactement égale à sa valeur actuelle (car $r=0$) :
\begin{equation*}
    \mathbb{E}[S_1] = S_0
\end{equation*}
\end{theorembox}

\begin{examplebox}[Calcul du prix via la probabilité risque-neutre]
\textbf{Étape 1 : Trouver la probabilité $p$}\\
Posons l'équation d'espérance de l'action :
\begin{equation*}
    110p + 90(1-p) = 100 \implies 20p + 90 = 100 \implies 20p = 10 \implies p = \frac{1}{2}
\end{equation*}

\textbf{Étape 2 : Évaluer l'option}\\
Dans ce monde risque-neutre, la valeur de l'option aujourd'hui est simplement l'espérance mathématique de ses flux futurs :
\begin{equation*}
    Opt_0 = \mathbb{E}[Flux] = p \times 10 + (1-p) \times 0 = \frac{1}{2}(10) + \frac{1}{2}(0) = 5
\end{equation*}
Nous retrouvons exactement le prix d'arbitrage calculé précédemment, mais avec un calcul beaucoup plus direct.
\end{examplebox}

\subsection{Généralisation et Théorie des Martingales}

\begin{examplebox}[Généralisation algébrique]
Considérons un produit dérivé générique payant $a + \beta$ dans l'état "Up" et $a$ dans l'état "Down".
\begin{itemize}[label=$\rightarrow$]
    \item \textbf{Par couverture} : On égalise les états pour trouver $\delta$ : $110\delta - (a + \beta) = 90\delta - a \implies \delta = \frac{\beta}{20}$. La valeur future garantie est $90(\frac{\beta}{20}) - a = 4,5\beta - a$. L'équation de non-arbitrage donne : $100(\frac{\beta}{20}) - Opt_0 = 4,5\beta - a \implies Opt_0 = a + 0,5\beta$.
    \item \textbf{Par risque-neutre} : L'espérance avec $p=0,5$ donne immédiatement $\mathbb{E}[Opt] = 0,5(a + \beta) + 0,5(a) = a + 0,5\beta$.
\end{itemize}
L'approche risque-neutre est considérablement plus simple à manipuler mathématiquement.
\end{examplebox}

\begin{theorembox}[Propriété de Martingale et Absence d'Arbitrage]
Pourquoi la tarification risque-neutre fonctionne-t-elle ? Une fois la probabilité $p$ fixée de manière à ce que $\mathbb{E}[S] = S_0$, l'opérateur espérance étant linéaire, cette relation s'applique à toute combinaison d'instruments sur le marché.
Si l'on note $B$ l'actif sans risque, on obtient le système suivant :
\begin{align*}
    \mathbb{E}[B_1] &= B_0 \\
    \mathbb{E}[S_1] &= S_0 \\
    \mathbb{E}[Opt_1] &= Opt_0
\end{align*}
La valeur actuelle de \textbf{tout} instrument est l'espérance de sa valeur future. Si un portefeuille vaut zéro aujourd'hui, son espérance future est zéro. Il est donc impossible que ce portefeuille soit toujours positif sans pouvoir être négatif : l'arbitrage est structurellement impossible sous cette probabilité.
\end{theorembox}

\begin{remarquebox}[Unicité du prix]
Les deux arguments (couverture et risque-neutre) fonctionnent en sens inverse :
\begin{itemize}[label=$\rightarrow$]
    \item L'argument risque-neutre montre qu'un prix précis est exempt d'arbitrage (il fixe une borne inférieure).
    \item L'argument par couverture montre que tous les autres prix peuvent être arbitrés (il fixe une borne supérieure).
\end{itemize}
Dans notre modèle binomial, ces deux bornes coïncident parfaitement. Le marché est dit "complet", et nous sommes laissés avec un prix d'arbitrage unique.
\end{remarquebox}

\subsection{La Tarification par Réplication}

\begin{intuitionbox}[La Loi du Prix Unique et la Réplication]
Il existe une troisième manière d'interpréter le prix d'arbitrage d'une option. Au lieu de vendre l'option et de nous couvrir, nous pouvons chercher à fabriquer un clone financier de cette option.
Si nous parvenons à construire un portefeuille composé d'actions et d'emprunts qui génère strictement les mêmes flux que l'option dans tous les états futurs possibles, le principe d'absence d'arbitrage impose que le coût de construction de ce portefeuille aujourd'hui soit exactement égal au prix de l'option.
\end{intuitionbox}

\begin{remarquebox}[L'argument géométrique]
Pourquoi est-il toujours possible de cloner l'option dans notre modèle ? 
La valeur d'un portefeuille d'actions et d'obligations à l'échéance est une fonction affine du prix de l'action, soit Valeur = Delta fois le prix de l'action + Montant obligataire. Graphiquement, c'est une ligne droite.
L'option définit deux points cibles à atteindre : son flux en cas de hausse, et son flux en cas de baisse. Or, par deux points, il passe une et une seule droite. Il existe donc une combinaison unique et exacte d'actions et d'emprunt qui réplique parfaitement l'option.
\end{remarquebox}

\begin{examplebox}[Réplication d'un Call sur Apple Inc.]
Imaginons que nous voulions évaluer une option d'achat (Call) sur l'action Apple (ticker : AAPL).
\begin{itemize}[label=$\rightarrow$]
    \item L'action AAPL cote aujourd'hui 200 \$.
    \item Demain, à la clôture, suite à une annonce de résultats, l'action sera soit à 220 \$ (Hausse), soit à 180 \$ (Baisse).
    \item Le taux d'intérêt à un jour d'un bon du Trésor américain (US Treasury Bill, ticker : BIL) est supposé nul (0\%) pour simplifier.
    \item Nous étudions un Call AAPL de prix d'exercice (Strike) à 200 \$.
\end{itemize}

Étape 1 : Identifier les flux cibles de l'option
\begin{itemize}[label=$\cdot$]
    \item Si AAPL monte à 220 \$, le Call vaut 20 \$.
    \item Si AAPL baisse à 180 \$, le Call vaut 0 \$.
\end{itemize}

Étape 2 : Construire le portefeuille de réplication
Nous cherchons à détenir une quantité Delta d'actions AAPL et un montant B en obligations US Treasury. Nous posons notre système d'équations pour le lendemain :
\begin{itemize}[label=$\cdot$]
    \item Scénario Hausse : 220 Delta + B = 20
    \item Scénario Baisse : 180 Delta + B = 0
\end{itemize}

En soustrayant la deuxième équation à la première, on obtient : 
40 Delta = 20, ce qui donne Delta = 0,5.
En remplaçant Delta dans la deuxième équation : 
180(0,5) + B = 0, ce qui donne 90 + B = 0, donc B = -90.

Étape 3 : Interprétation financière et tarification
Le calcul mathématique nous dicte la stratégie exacte à suivre aujourd'hui :
\begin{itemize}[label=$\cdot$]
    \item Acheter 0,5 action AAPL.
    \item Vendre à découvert pour 90 \$ de bons du Trésor américain (ce qui équivaut à emprunter 90 \$ à la banque).
\end{itemize}

Vérifions ce clone demain : si AAPL vaut 180 \$, notre demi-action vaut 90 \$, ce qui sert tout juste à rembourser notre dette de 90 \$ au Trésor. Il nous reste 0 \$ (exactement comme le Call). Si AAPL vaut 220 \$, notre demi-action vaut 110 \$, nous remboursons les 90 \$, il nous reste 20 \$ net en poche (exactement comme le Call).

Quel est le coût de création de ce clone aujourd'hui ?
Coût net = Achat des actions AAPL - Argent emprunté via les Treasury Bills
Coût net = 0,5 fois 200 \$ - 90 \$ = 100 \$ - 90 \$ = 10 \$.

Le clone parfait coûte 10 \$. Par conséquent, le prix d'arbitrage de l'option Call AAPL sur le marché doit être très exactement de 10 \$.
\end{examplebox}

\begin{theorembox}[Symétrie de la réplication]
Sur un marché complet, l'action, l'obligation sans risque et l'option sont structurellement liés. L'équation de valorisation peut être manipulée algébriquement pour répliquer n'importe quel instrument à partir des deux autres :
\begin{itemize}[label=$\rightarrow$]
    \item Hedging : Action possédée + Option vendue permet de générer des flux identiques à une Obligation sans risque.
    \item Réplication : Action possédée + Obligation empruntée permet de générer des flux identiques à une Option d'achat.
    \item Synthèse d'actif : Option achetée + Obligation détenue permet de générer des flux identiques à une Action.
\end{itemize}
\end{theorembox}

\subsection{Les Limites du Modèle : Le Marché Incomplet}

\begin{intuitionbox}[L'ajout d'un troisième état]
Le modèle binomial à deux états est parfait mathématiquement, mais irréaliste économiquement. Pour s'approcher de la réalité, supposons un instant que l'action (actuellement à 100) puisse prendre trois valeurs demain : monter à 110 (A), stagner à 100 (C), ou baisser à 90 (B). Le Call a toujours un Strike de 100.
Essayons d'appliquer nos méthodes de valorisation à cet univers à trois états.
\end{intuitionbox}

\begin{proofbox}[L'échec de la couverture et l'apparition d'inégalités]
Tentons de nous couvrir en vendant l'option et en achetant une quantité y d'actions. Notre portefeuille demain vaudra :
\begin{itemize}[label=$\cdot$]
    \item État A (110) : $110y - 10$
    \item État C (100) : $100y$
    \item État B (90) : $90y$
\end{itemize}
Pour éliminer le risque, ces trois montants doivent être égaux. Mathématiquement, c'est impossible : nous avons 3 équations simultanées à résoudre mais une seule variable (y).
Si nous utilisons le ratio y = 0,5 du modèle précédent, notre portefeuille vaudra 45 (en A), 50 (en C) ou 45 (en B). Le risque n'est pas éliminé, mais nous savons que le portefeuille vaudra au minimum 45. Par conséquent, sa valeur initiale doit être d'au moins 45 :
\begin{equation*}
    100(0,5) - Opt_0 \ge 45 \implies Opt_0 \le 5
\end{equation*}
L'argument d'arbitrage ne donne plus un prix exact, mais une borne supérieure.
\end{proofbox}

\begin{examplebox}[L'échec de la probabilité risque-neutre unique]
Cherchons une probabilité risque-neutre telle que l'espérance de l'action demain soit 100.
Soit $P_A, P_B, P_C$ les probabilités des trois états.
\begin{equation*}
    110 P_A + 100 P_C + 90 P_B = 100
\end{equation*}
Sachant que $P_C = 1 - P_A - P_B$, l'équation se simplifie pour donner : $P_A = P_B$.
Posons $P_A = P_B = p$. La probabilité $p$ peut être n'importe quel chiffre compris entre 0 et 0,5. Il existe donc une infinité de probabilités risque-neutres valides.
Le prix de l'option étant son espérance, nous avons :
\begin{equation*}
    Opt_0 = 10 \times p + 0 \times P_C + 0 \times P_B = 10p
\end{equation*}
Puisque p est compris entre 0 et 0,5, le prix du Call est un intervalle compris entre 0 et 5.
\end{examplebox}

\begin{theorembox}[Définition d'un Marché Incomplet]
Le modèle à trois états illustre le concept fondamental de marché incomplet. Un marché est incomplet lorsqu'il y a plus d'états futurs du monde que d'actifs financiers indépendants permettant de s'y couvrir.
Dans un tel marché :
\begin{itemize}[label=$\rightarrow$]
    \item La réplication parfaite d'un produit dérivé est impossible.
    \item L'absence d'opportunité d'arbitrage ne suffit plus à fixer un prix unique.
    \item Les mathématiques ne fournissent qu'un intervalle de prix exempts d'arbitrage (ici [0, 5]).
\end{itemize}
Le prix exact qui sera négocié sur le marché ne dépendra plus d'une équation, mais des préférences au risque et de la psychologie des acteurs du marché.
\end{theorembox}

\subsection{Modèle à Plusieurs Périodes}

\begin{intuitionbox}[L'astuce du temps continu pour sauver le modèle]
Nous venons de voir qu'ajouter un troisième état final (stagnation à 100) détruit l'unicité du prix et rend le marché incomplet. 
Pour modéliser des mouvements de prix réalistes sans perdre la puissance mathématique de l'arbitrage, la solution consiste à découper le temps. Si nous supposons que les prix évoluent de manière continue, l'action ne fait pas un grand saut en une journée, elle fait plusieurs petits sauts successifs.
\end{intuitionbox}

\newpage

\begin{examplebox}[Construction d'un arbre binomial à deux pas]
Plutôt que d'envisager une variation de plus ou moins 10 sur une journée entière, nous découpons la journée en deux demi-journées, avec des variations de plus ou moins 5.
\begin{itemize}[label=$\rightarrow$]
    \item Instant initial (t = 0) : L'action vaut 100.
    \item Étape intermédiaire (t = 0,5) : L'action fait face à un univers à deux états. Elle monte à 105 (Nœud D) ou baisse à 95 (Nœud E).
    \item Étape finale (t = 1) : À partir de chaque nœud intermédiaire, l'action fait de nouveau face à deux états.
    \begin{itemize}[label=$\cdot$]
        \item Depuis 105 : elle monte à 110 (État A) ou baisse à 100 (État C).
        \item Depuis 95 : elle monte à 100 (État C) ou baisse à 90 (État B).
    \end{itemize}
\end{itemize}
L'état C (100) est atteint de deux manières différentes : une hausse suivie d'une baisse, ou une baisse suivie d'une hausse. C'est ce que l'on appelle un arbre recombinant.
\end{examplebox}

\vspace{0.5cm}
\begin{center}
\begin{tikzpicture}[>=stealth, sloped]
    % Définition des styles pour les noeuds
    \tikzstyle{node_style} = [circle, draw=vscodeBlue, fill=white, thick, minimum size=1cm, inner sep=2pt, font=\small]
    \tikzstyle{time_style} = [font=\small\color{gray}]

    % Les axes de temps (Placés plus haut à y=4.5 pour éviter le chevauchement)
    \node[time_style] at (0, 4.5) {Aujourd'hui ($t=0$)};
    \node[time_style] at (4, 4.5) {Mi-journée ($t=0,5$)};
    \node[time_style] at (8, 4.5) {Demain ($t=1$)};

    % Placement des noeuds
    % t=0
    \node[node_style] (S0) at (0, 0) {100};

    % t=0.5
    \node[node_style] (D) at (4, 1.5) {105};
    \node[node_style] (E) at (4, -1.5) {95};

    % t=1 (Le noeud A est à y=3)
    \node[node_style] (A) at (8, 3) {110};
    \node[node_style] (C) at (8, 0) {100};
    \node[node_style] (B) at (8, -3) {90};

    % Dessin des flèches
    \draw[->, thick, vscodeBlue] (S0) -- (D) node[midway, above] {Hausse};
    \draw[->, thick, vscodeBlue] (S0) -- (E) node[midway, below] {Baisse};

    \draw[->, thick, vscodeBlue] (D) -- (A);
    \draw[->, thick, vscodeBlue] (D) -- (C);
    
    \draw[->, thick, vscodeBlue] (E) -- (C);
    \draw[->, thick, vscodeBlue] (E) -- (B);
    
    % Noms des états à côté des noeuds finaux
    \node[right=0.2cm of A, font=\small] {État A};
    \node[right=0.2cm of C, font=\small] {État C};
    \node[right=0.2cm of B, font=\small] {État B};
    
    % Noms des états intermédiaires
    \node[above=0.2cm of D, font=\small] {Nœud D};
    \node[below=0.2cm of E, font=\small] {Nœud E};

\end{tikzpicture}
\end{center}
\vspace{0.5cm}

\begin{remarquebox}[Le retour à la complétude du marché]
Ce découpage temporel est une avancée théorique majeure. En regardant l'arbre globalement à t=1, nous avons bien trois états finaux possibles (110, 100, 90). Cependant, si nous nous plaçons à n'importe quel nœud spécifique de l'arbre, il n'y a que deux chemins possibles pour l'étape suivante. 
Le marché est donc localement complet à chaque nœud. Nous pourrons utiliser nos techniques d'évaluation (Couverture ou Risque-Neutre) à chaque intersection pour déduire un prix exact, en travaillant à rebours depuis l'échéance finale jusqu'à aujourd'hui.
\end{remarquebox}

\subsection{Évaluation par Couverture dans un Arbre à Deux Périodes}

\begin{intuitionbox}[Le principe d'induction rétrograde]
Pour calculer la valeur de l'option aujourd'hui, nous devons procéder à l'envers (calcul à rebours). Puisque le marché est complet à chaque nœud, nous pouvons appliquer notre argument de couverture locale. Nous calculons d'abord la valeur de l'option à la mi-journée (nœuds D et E) en utilisant les flux finaux (nœuds A, B et C). Une fois ces valeurs intermédiaires connues, nous recommençons l'opération pour trouver la valeur initiale à t=0.
\end{intuitionbox}

\newpage

\begin{examplebox}[Calcul pas à pas par couverture (Delta Hedging)]
Considérons notre option d'achat (Call) de prix d'exercice K = 100. Son profil de flux à l'échéance (t=1) est : 10 en A, 0 en C, 0 en B. Le taux d'intérêt sans risque est toujours supposé nul.

\begin{enumerate}[label=Étape \arabic* :, leftmargin=*]
    
    \item Évaluation au nœud E (L'action vaut 95 à t=0,5)
    \begin{itemize}[label=$\cdot$]
        \item Si nous sommes au nœud E, l'action finira à l'échéance soit à 100 (État C), soit à 90 (État B). 
        \item Dans les deux cas, le Call de Strike 100 ne vaudra strictement rien (0).
        \item Conclusion immédiate : La valeur de l'option au nœud E est Opt(E) = 0.
    \end{itemize}
    
    \vspace{0.2cm}
    \item Évaluation au nœud D (L'action vaut 105 à t=0,5)
    \begin{itemize}[label=$\cdot$]
        \item À partir de D, l'action peut monter à 110 (le Call vaudra 10) ou baisser à 100 (le Call vaudra 0).
        \item Construisons le portefeuille de couverture : nous achetons une quantité delta d'actions.
        \item Valeur future du portefeuille : 110 delta - 10 (si hausse) et 100 delta - 0 (si baisse).
        \item Pour éliminer le risque, nous égalisons : 110 delta - 10 = 100 delta, ce qui donne 10 delta = 10, soit delta = 1.
        \item Avec 1 action, la valeur garantie à t=1 est de 100.
        \item Par arbitrage, le portefeuille au nœud D vaut 100 : 105(1) - Opt(D) = 100, ce qui implique Opt(D) = 5.
    \end{itemize}

    \vspace{0.2cm}
    \item Évaluation au nœud initial (L'action vaut 100 à t=0)
    \begin{itemize}[label=$\cdot$]
        \item Nous connaissons désormais les valeurs à la mi-journée : 5 si l'action monte à 105, 0 si elle baisse à 95.
        \item Construisons le portefeuille initial : nous achetons une nouvelle quantité delta d'actions.
        \item Valeur future (à t=0,5) : 105 delta - 5 (si hausse vers D) et 95 delta - 0 (si baisse vers E).
        \item Égalisation des flux : 105 delta - 5 = 95 delta, ce qui donne 10 delta = 5, soit delta = 0,5.
        \item Avec 0,5 action, la valeur garantie à t=0,5 est de 95(0,5) = 47,5.
        \item Par arbitrage, la valeur initiale du portefeuille aujourd'hui doit être de 47,5 : 100(0,5) - Opt(0) = 47,5.
        \item Le prix final de l'option aujourd'hui est Opt(0) = 2,5.
    \end{itemize}

\end{enumerate}
\end{examplebox}

\begin{remarquebox}[Couverture dynamique]
Cet exemple démontre mathématiquement la notion de couverture dynamique évoquée précédemment. Remarquez que le delta (la quantité d'actions à détenir pour se couvrir) n'est pas constant. Il vaut 0,5 au début de la journée. Si l'action monte à 105, le vendeur de l'option devra ajuster son portefeuille en achetant 0,5 action supplémentaire pour que son delta passe à 1. La couverture parfaite exige un ajustement à chaque pas de temps.
\end{remarquebox}

\subsection{Évaluation Risque-Neutre dans un Arbre à Deux Périodes}

\begin{intuitionbox}[L'avantage de la méthode probabiliste]
Bien que l'argument de couverture nous garantisse un prix mathématiquement parfait, les calculs de portefeuilles deviennent vite fastidieux. L'approche par la probabilité risque-neutre, bien que plus abstraite, est beaucoup plus simple à appliquer. Elle donnera toujours un résultat strictement identique à la méthode par couverture.
Il existe deux manières de l'utiliser : en remontant l'arbre pas à pas (induction rétrograde), ou en calculant directement les probabilités globales d'atteindre chaque état final.
\end{intuitionbox}

\newpage

\begin{examplebox}[Méthode 1 : Calcul pas à pas (Induction rétrograde)]
Nous reprenons le même arbre. Le taux d'intérêt est nul, donc l'espérance du prix futur de l'action doit égaler son prix actuel à chaque nœud.

\begin{enumerate}[label=Étape \arabic* :, leftmargin=*]
    \item Évaluation au nœud D (L'action vaut 105)
    \begin{itemize}[label=$\cdot$]
        \item L'action peut passer à 110 ou 100. La probabilité p (Hausse) doit vérifier : 110p + 100(1-p) = 105.
        \item La seule probabilité valide est p = 1/2.
        \item La valeur de l'option est son espérance : Opt(D) = 0,5(10) + 0,5(0) = 5.
    \end{itemize}

    \vspace{0.2cm}
    \item Évaluation au nœud E (L'action vaut 95)
    \begin{itemize}[label=$\cdot$]
        \item Les flux finaux de l'option sont nuls dans tous les scénarios (États C et B).
        \item L'espérance est de zéro quelle que soit la probabilité. Opt(E) = 0.
    \end{itemize}

    \vspace{0.2cm}
    \item Évaluation initiale (L'action vaut 100)
    \begin{itemize}[label=$\cdot$]
        \item L'action passe à 105 ou 95. Pour que l'espérance soit 100, la probabilité de Hausse doit de nouveau être p = 1/2.
        \item La valeur initiale de l'option est l'espérance de ses valeurs futures intermédiaires (Opt(D) et Opt(E)).
        \item Opt(0) = 0,5(5) + 0,5(0) = 2,5.
    \end{itemize}
\end{enumerate}
Le résultat est identique au prix d'arbitrage obtenu par la méthode de couverture en delta.
\end{examplebox}

\begin{examplebox}[Méthode 2 : Calcul par les probabilités globales des chemins]
Plutôt que de calculer les valeurs intermédiaires de l'option, nous pouvons calculer directement la probabilité risque-neutre de nous retrouver dans chaque état final à t=1, puis calculer l'espérance finale en une seule fois.
Sachant qu'à chaque nœud la probabilité de Hausse est de 1/2 :

\begin{itemize}[label=$\cdot$]
    \item État A (Action à 110) : Il faut 2 hausses successives. Probabilité = 1/2 fois 1/2 = 1/4.
    \item État C (Action à 100) : Il faut 1 hausse et 1 baisse (ou inversement). Probabilité = 1/4 + 1/4 = 1/2.
    \item État B (Action à 90) : Il faut 2 baisses successives. Probabilité = 1/2 fois 1/2 = 1/4.
\end{itemize}

Calcul direct de l'espérance de l'option aujourd'hui :
Opt(0) = 1/4(10) + 1/2(0) + 1/4(0) = 10/4 = 2,5.

Cette méthode est extrêmement puissante car elle permet de sauter toutes les étapes intermédiaires de calcul.
\end{examplebox}

\begin{theorembox}[Généralisation et Tarification Dynamique]
L'unique caractéristique qualitative cruciale de cet argument est qu'à chaque pas de temps, l'actif ne peut évoluer que vers deux nouvelles valeurs possibles. C'est ce qui garantit la complétude locale du marché.
\begin{itemize}[label=$\rightarrow$]
    \item Le nombre total de pas de temps n'a aucune importance mathématique. On peut subdiviser le temps en autant d'étapes que l'on souhaite (10, 100, 1000 étapes). Tant que l'arbre reste binomial à chaque nœud, le prix d'arbitrage sera unique.
    \item L'interprétation par la réplication reste valide : le prix de l'option (2,5) est exactement la somme d'argent qu'il faut investir aujourd'hui pour fabriquer un clone parfait de l'option.
    \item Ce clone nécessite un rééquilibrage continu. La quantité d'actions à détenir (le Delta) et le montant de cash (emprunt ou placement) changent à chaque nœud de l'arbre en fonction du temps et du prix de l'action. C'est l'essence même de la gestion de portefeuille dynamique.
\end{itemize}
\end{theorembox}

\subsection{Généralisation à N Périodes}

\begin{intuitionbox}[La mécanique de l'arbre global]
Il n'y a rien de magique dans un modèle à deux étapes. Tant que chaque nœud de l'arbre se divise exactement en deux branches, la logique d'induction rétrograde garantit un prix sans arbitrage unique, que l'arbre compte 2, 4 ou 10 000 étapes.
Au lieu de calculer à rebours (nœud par nœud, couche par couche) depuis l'échéance finale jusqu'à aujourd'hui, nous pouvons directement calculer les probabilités globales d'atteindre chaque nœud final de l'arbre, puis faire l'espérance mathématique des flux finaux.
\end{intuitionbox}

\newpage

\begin{examplebox}[Exemple d'un arbre à 4 pas de temps (Taux d'intérêt = 0)]
Imaginons un arbre à 4 étapes, où le cours de l'action monte ou baisse de 2,5 à chaque pas. L'action part de 100 à t=0. 
Puisque le taux d'intérêt est nul, la probabilité risque-neutre de hausse (p) est obligatoirement de 0,5 à chaque nœud pour maintenir l'espérance locale.
Calculons la probabilité risque-neutre globale de chaque état final à la 4ème étape :

\begin{itemize}[label=$\cdot$]
    \item Arrivée à 110 : Il faut exactement 4 hausses successives.
    Il n'y a qu'un seul chemin. Probabilité : $1 \times (1/2)^4 = 1/16$.
    
    \item Arrivée à 105 : Il faut exactement 3 hausses et 1 baisse.
    Le nombre de combinaisons possibles est donné par le coefficient binomial $C_4^1 = 4$.
    Probabilité : $4 \times (1/2)^4 = 4/16 = 1/4$.
    
    \item Arrivée à 100 : Il faut exactement 2 hausses et 2 baisses.
    Nombre de combinaisons $C_4^2 = 6$.
    Probabilité : $6 \times (1/2)^4 = 6/16 = 3/8$.
    
    \item Arrivée à 95 : Il faut exactement 1 hausse et 3 baisses.
    Nombre de combinaisons $C_4^3 = 4$.
    Probabilité : $4 \times (1/2)^4 = 4/16 = 1/4$.
    
    \item Arrivée à 90 : Il faut exactement 0 hausse et 4 baisses.
    Il n'y a qu'un seul chemin. Probabilité : $1 \times (1/2)^4 = 1/16$.
\end{itemize}

Tarification d'un Call de Strike 100 :
Nous multiplions la probabilité de chaque état par la valeur intrinsèque du Call dans cet état.
Opt(K=100) = (1/16) fois (110 - 100) + (1/4) fois (105 - 100) + 0 + 0 + 0
Opt(K=100) = 10/16 + 5/4 = 10/16 + 20/16 = 30/16 = 15/8.
\end{examplebox}

\vspace{0.5cm}
\begin{center}
\begin{tikzpicture}[>=stealth, sloped]
    % Style pour les mini noeuds d'un grand arbre
    \tikzstyle{mini_node} = [circle, draw=vscodeBlue, fill=codeBackground, thick, minimum size=0.4cm, inner sep=0pt, font=\tiny]
    \tikzstyle{time_style} = [font=\small\color{gray}]

    % Dimensions de l'arbre
    \def\N{5} % Nombre d'étapes
    \def\dx{1.8} % Espace horizontal entre les étapes
    \def\dy{0.8} % Espace vertical entre les noeuds

    % Lignes de temps
    \foreach \i in {0,...,\N} {
        \node[time_style] at (\i*\dx, \N*\dy*0.5 + 0.8) {$t=\i$};
    }

    % Génération automatique des noeuds et des liens (Arbre Recombinant)
    \foreach \i in {0,...,\N} {
        \foreach \j in {0,...,\i} {
            % Calcul de la position Y : Centré autour de Y=0
            \pgfmathsetmacro{\ypos}{(\i - 2*\j)*\dy*0.5}
            \node[mini_node] (N-\i-\j) at (\i*\dx, \ypos) {};
            
            % Dessiner les liens vers la génération précédente (si on n'est pas à i=0)
            \ifnum\i>0
                \pgfmathtruncatemacro{\prevI}{\i-1} % CORRECTION ICI
                % Lien depuis le noeud supérieur (baisse)
                \ifnum\j>0
                    \pgfmathtruncatemacro{\prevJup}{\j-1} % CORRECTION ICI
                    \draw[->, vscodeBlue!60] (N-\prevI-\prevJup) -- (N-\i-\j);
                \fi
                % Lien depuis le noeud inférieur (hausse)
                \ifnum\j<\i
                    \pgfmathtruncatemacro{\prevJdown}{\j} % CORRECTION ICI
                    \draw[->, vscodeBlue!60] (N-\prevI-\prevJdown) -- (N-\i-\j);
                \fi
            \fi
        }
    }
    
    % Étiquette pour montrer la multiplication des chemins
    \node[right=0.2cm of N-\N-0, font=\scriptsize] {5 Hausses};
    \node[right=0.2cm of N-\N-5, font=\scriptsize] {5 Baisses};
    \node[right=0.2cm of N-\N-2, font=\scriptsize] {Chemins multiples};
    
\end{tikzpicture}
\end{center}
\vspace{0.5cm}

\begin{theorembox}[La Formule Globale pour N pas de temps]
Considérons une action de prix initial S. L'arbre compte $N$ étapes. À chaque étape, l'action monte ou baisse d'un montant absolu $x$.
À la fin de l'arbre, la valeur de l'action dépend uniquement du nombre total de hausses (noté $j$). Si l'action a fait $j$ hausses, elle a obligatoirement fait $(N - j)$ baisses.
Le prix final de l'action $S_j$ dans l'état avec $j$ hausses s'écrit :
\begin{equation*}
    S_j = S + j(x) - (N - j)(x) = S + (2j - N)x
\end{equation*}

Puisque la probabilité d'une hausse est de $1/2$ (avec un taux nul), la probabilité totale d'atteindre ce prix final $S_j$ suit une distribution binomiale :
\begin{equation*}
    Prob(S_j) = \frac{1}{2^N} \binom{N}{j}
\end{equation*}
Où $\binom{N}{j}$ représente le nombre de chemins différents permettant d'obtenir $j$ hausses sur $N$ lancers.

Par conséquent, la valeur actuelle d'un produit dérivé qui génère un flux $f(S)$ à l'échéance est simplement l'espérance de tous les flux terminaux pondérés par cette probabilité :
\begin{equation*}
    Opt_0 = \frac{1}{2^N} \sum_{j=0}^{N} \binom{N}{j} f\left(S + (2j - N)x\right)
\end{equation*}
\end{theorembox}

\newpage

\begin{intuitionbox}[Vers le temps continu et Black-Scholes]
L'objectif ultime de cette modélisation mathématique n'est pas de calculer un arbre géant à la main, mais de comprendre ce qui se passe mathématiquement lorsque l'on fait tendre le nombre d'étapes $N$ vers l'infini. 
Pour que la limite ait un sens mathématique et économique, il faudra réduire la taille de la variation de prix ($x$) au fur et à mesure que le nombre d'étapes ($N$) augmente, tout en préservant certaines propriétés de la volatilité de l'action. 
C'est précisément ce passage à la limite qui nous fera basculer de l'arbre binomial discret au mouvement brownien continu, fondement du modèle de Black-Scholes.
\end{intuitionbox}